\documentclass{beamer}

\usetheme{Berlin}
\usecolortheme{Orchid}
\useoutertheme{miniframes}
\setbeamertemplate{navigation symbols}{}

\usepackage[utf8]{inputenc}
\usepackage[T1]{fontenc}
\usepackage{amsmath}
\usepackage{amssymb}
\usepackage{booktabs}
\usepackage{graphicx}
\graphicspath{{../images/}}

\usepackage{xcolor}
\usepackage{listings}
\usepackage{hyperref}
\usepackage{tikz}
\usetikzlibrary{arrows.meta,positioning,shapes.geometric,calc}

\lstset{
  language=Java,
  basicstyle=\ttfamily\small,
  keywordstyle=\color{blue},
  commentstyle=\color{gray},
  stringstyle=\color{teal},
  showstringspaces=false,
  columns=fullflexible,
  breaklines=true
}

% Per-slide source line (references shown on the slide itself).
\newcommand{\slidesources}[1]{%
  \vfill
  {\tiny\color{gray}\raggedright\sloppy\parbox{\linewidth}{Sources: #1}\par}%
}

% Unit-wise source strings for this subject (used by slides/notes).
% Path is relative to the lecture `latex/` folder (the compilation CWD).
% OOPS with Java (CSEG1044) — unit-wise sources
%
% These are short, student-facing reference strings intended to be used with:
%   - \slidesources{...}  (in beamer slides)
%   - \notesources{...}   (in lecture notes)
%
% Style inspiration: other subjects in My-Subjects/ use semicolon-separated sources
% and include an "accessed" date for web documentation.

\newcommand{\OOPSUnitISources}{Schildt, \textit{Java: A Beginner's Guide} (10e), McGraw-Hill (Java fundamentals + OOP basics).; Oracle Java Tutorials: Getting Started + Language Basics + Classes and Objects (accessed \today).; Java SE API docs (e.g., \texttt{String}, \texttt{Scanner}) (accessed \today).}

\newcommand{\OOPSUnitIISources}{Schildt, \textit{Java: The Complete Reference} (12e), McGraw-Hill (classes, inheritance, packages, interfaces).; Bloch, \textit{Effective Java} (3e), Addison-Wesley, 2018 (design choices; interfaces vs abstract classes).; Oracle Java Tutorials: Classes and Objects + Interfaces + Packages (accessed \today).; Java SE API docs (e.g., \texttt{Comparable}, \texttt{Runnable}) (accessed \today).}

\newcommand{\OOPSUnitIIISources}{Schildt, \textit{Java: The Complete Reference} (12e) (nested/inner classes, exceptions, threads, I/O).; Oracle Java Tutorials: Nested Classes + Exceptions + Concurrency + Basic I/O (accessed \today).; Java SE API docs (e.g., \texttt{Thread}, \texttt{AutoCloseable}) (accessed \today).}

\newcommand{\OOPSUnitIVSources}{Schildt, \textit{Java: The Complete Reference} (12e) (generics, lambdas, Swing, JDBC).; Bloch, \textit{Effective Java} (3e) (generics + lambdas).; Oracle Java Tutorials: Generics + Lambda Expressions + Swing + JDBC Basics (accessed \today).; Java SE API docs (e.g., \texttt{java.util.function}, \texttt{javax.swing}, \texttt{java.sql}) (accessed \today).}

\newcommand{\OOPSUnitVSources}{Schildt, \textit{Java: The Complete Reference} (12e) (Collections Framework + wrapper classes).; Oracle Java docs: Collections Framework overview (accessed \today).; Java SE API docs (\texttt{java.util} collections; wrapper classes in \texttt{java.lang}) (accessed \today).}

\newcommand{\OOPSUnitVISources}{Git SCM: \textit{Pro Git} (2e) (version control workflow), accessed \today.; GitHub Docs (repositories, issues, pull requests), accessed \today.; Oracle Java Tutorials: Swing + JDBC (accessed \today).; JUnit 5 User Guide (accessed \today).; Apache Maven Project: Getting Started (accessed \today).; Gradle User Manual: Getting Started (accessed \today).; UML class diagrams: UML 2.x overview/spec (accessed \today).}

\newcommand{\OOPSMiscWhyInterfacesSources}{Schildt, \textit{Java: The Complete Reference} (12e) (interfaces).; Bloch, \textit{Effective Java} (3e) (prefer interfaces where appropriate).; Oracle Java Tutorials: Interfaces (accessed \today).; Java SE API docs (e.g., \texttt{Comparable}, \texttt{Runnable}) (accessed \today).}



\title[OOPS with Java]{Object Oriented Programming (CSEG1044)}
\subtitle{Unit 03 -- Lecture 04: Java I/O Streams + FileReader/FileWriter}
\author{Tofik Ali}
\institute{School of Computer Science, UPES Dehradun}
\date{\today}

\begin{document}

\begin{frame}[fragile]
  \titlepage
  \vspace{0.5em}
  \slidesources{\OOPSUnitIIISources}
\end{frame}

\begin{frame}{Quick Links}
  \centering
  \hyperlink{sec:overview}{\beamerbutton{Overview}}\hspace{0.6em}
  \hyperlink{sec:concepts}{\beamerbutton{Key Topics}}\hspace{0.6em}
  \hyperlink{sec:demo}{\beamerbutton{Demo}}\hspace{0.6em}
  \hyperlink{sec:summary}{\beamerbutton{Summary}}
  \slidesources{\OOPSUnitIIISources}
\end{frame}

\begin{frame}{Agenda}
  \tableofcontents
  \slidesources{\OOPSUnitIIISources}
\end{frame}

\section{Overview}
\label{sec:overview}

\begin{frame}{Learning Outcomes}
  \begin{itemize}[<+->]
    \item Explain stream model; choose byte vs character streams.
    \item Read/write text files with \texttt{FileReader}/\texttt{FileWriter} and buffering.
    \item Use \texttt{try-with-resources} for exception-safe I/O code.
  \end{itemize}
  \slidesources{\OOPSUnitIIISources}
\end{frame}

\section{Key Topics}
\label{sec:concepts}

\begin{frame}{Unit 03: Nested Classes, Exceptions, Multithreading \& IO Streams}
  \begin{itemize}[<+->]
    \item Lecture focus: Java I/O Streams + FileReader/FileWriter
  \end{itemize}
  \slidesources{\OOPSUnitIIISources}
\end{frame}

\begin{frame}{Today: Roadmap}
  \begin{itemize}[<+->]
    \item Stream model; byte streams vs character streams
    \item File I/O with `FileReader` / `FileWriter` (+ buffering)
    \item Exception-safe code via `try-with-resources` and good error messages
  \end{itemize}
  \slidesources{\OOPSUnitIIISources}
\end{frame}

\begin{frame}{Stream model (big picture)}
  \begin{itemize}[<+->]
    \item Stream = sequence of data flowing from source to destination.
    \item Input stream reads; output stream writes.
    \item Streams can be chained (file stream wrapped with buffer).
  \end{itemize}
  \slidesources{\OOPSUnitIIISources}
\end{frame}

\begin{frame}{Byte vs character streams}
  \begin{itemize}[<+->]
    \item Byte streams: \texttt{InputStream}/\texttt{OutputStream} (binary: images, PDFs).
    \item Character streams: \texttt{Reader}/\texttt{Writer} (text files).
  \end{itemize}
  \begin{block}{Rule of thumb}
    Use \texttt{Reader/Writer} for text; use \texttt{InputStream/OutputStream} for binary.
  \end{block}
  \slidesources{\OOPSUnitIIISources}
\end{frame}

\begin{frame}{Text file I/O: \texttt{FileReader}/\texttt{FileWriter}}
  \begin{itemize}[<+->]
    \item \texttt{FileReader}: reads characters from a file.
    \item \texttt{FileWriter}: writes characters to a file.
    \item Wrap with buffering for performance and convenience.
  \end{itemize}
  \slidesources{\OOPSUnitIIISources}
\end{frame}

\begin{frame}{Buffering: why it matters}
  \begin{itemize}[<+->]
    \item Disk I/O is slow; reading/writing character-by-character is inefficient.
    \item \texttt{BufferedReader} reads chunks and gives \texttt{readLine()}.
    \item \texttt{BufferedWriter} writes efficiently and provides \texttt{newLine()}.
  \end{itemize}
  \slidesources{\OOPSUnitIIISources}
\end{frame}

\begin{frame}{\texttt{try-with-resources} (safe cleanup)}
  \begin{itemize}[<+->]
    \item Ensures files are closed even if an exception occurs.
    \item Syntax: \texttt{try (BufferedReader br = ...)\{...\}}
    \item Cleaner than manual \texttt{close()} in \texttt{finally}.
  \end{itemize}
  \slidesources{\OOPSUnitIIISources}
\end{frame}

\begin{frame}{Common beginner issue: file not found}
  \begin{itemize}[<+->]
    \item Paths are relative to the working directory (where you run \texttt{java}).
    \item Check spelling, extension, and location.
    \item Use absolute path if needed during debugging.
  \end{itemize}
  \slidesources{\OOPSUnitIIISources}
\end{frame}

\begin{frame}[fragile]{Working directory + Windows path strings}
  \begin{itemize}[<+->]
    \item Print working directory when debugging:
      \begin{lstlisting}
System.out.println(new java.io.File(".").getAbsolutePath());
      \end{lstlisting}
    \item Windows paths in Java strings need escaping:
      \begin{itemize}
        \item wrong: \verb|"D:\data\marks.txt"|
        \item correct: \verb|"D:\\data\\marks.txt"| or \verb|"D:/data/marks.txt"|
      \end{itemize}
  \end{itemize}
  \slidesources{\OOPSUnitIIISources}
\end{frame}


\section{Demo / Example}
\label{sec:demo}

\begin{frame}[fragile]{Example (1/2): read + compute (safe I/O)}
\begin{lstlisting}
import java.io.*;

public class MarksReport {
  public static void main(String[] args) {
    try (BufferedReader br = new BufferedReader(new FileReader("marks.txt"));
         BufferedWriter bw = new BufferedWriter(new FileWriter("report.txt"))) {

      int sum = 0, max = Integer.MIN_VALUE, count = 0;
      String line;
      while ((line = br.readLine()) != null) {
        line = line.trim();
        if (line.isEmpty()) continue;
        int x = Integer.parseInt(line);
        sum += x; if (x > max) max = x; count++;
      }
\end{lstlisting}
  \slidesources{\OOPSUnitIIISources}
\end{frame}

\begin{frame}[fragile]{Example (2/2): write report + handle errors}
\begin{lstlisting}
      double avg = (count == 0) ? 0.0 : (sum / (double) count);
      bw.write("Count=" + count); bw.newLine();
      bw.write("Max=" + max);     bw.newLine();
      bw.write(String.format("Avg=%.2f", avg)); bw.newLine();
    }
    catch (IOException e) {
      System.out.println("I/O error: " + e.getMessage());
    }
    catch (NumberFormatException e) {
      System.out.println("Bad number in file: " + e.getMessage());
    }
  }
}
\end{lstlisting}
  \slidesources{\OOPSUnitIIISources}
\end{frame}

\begin{frame}{Mini Exercises (2--3 mins)}
  \begin{enumerate}[<+->]
    \item Text file: byte stream or character stream?
    \item Why is buffering useful?
    \item What does \texttt{try-with-resources} guarantee?
  \end{enumerate}
  \slidesources{\OOPSUnitIIISources}
\end{frame}


\section{Summary}
\label{sec:summary}

\begin{frame}{Key Takeaways}
  \begin{itemize}[<+->]
    \item Streams model I/O; choose byte vs character based on data type.
    \item Use buffering for efficiency and line-based reading.
    \item Use \texttt{try-with-resources} for safe closing and cleaner code.
  \end{itemize}
  \slidesources{\OOPSUnitIIISources}
\end{frame}

\begin{frame}{Checkpoint Questions}
  \begin{enumerate}[<+->]
    \item When do you choose \texttt{InputStream} over \texttt{Reader}?
    \item Why do we prefer \texttt{BufferedReader} for reading lines?
    \item What is the working directory and why does it matter for file paths?
  \end{enumerate}
  \slidesources{\OOPSUnitIIISources}
\end{frame}

\begin{frame}{Next Unit Preview}
  \textbf{Next:} Unit 04 -- Generics Fundamentals (Generic Classes)
  \vspace{0.6em}
  \begin{itemize}[<+->]
    \item We will learn type-safe reusable code using \texttt{<T>}.
  \end{itemize}
  \slidesources{\OOPSUnitIIISources}
\end{frame}

\end{document}
